\section{Ethics, safety and governance}

Embodied medical intelligence changes the object of governance. A conventional clinical model primarily transforms information into a recommendation; an embodied system can also sense continuously, move through shared spaces, manipulate tissue or equipment, and deliver an intervention. Its risk therefore emerges from a coupled perception--decision--action loop rather than from prediction accuracy alone. Ambient sensing can support earlier and more contextual care, but it can also turn homes, wards and bodies into persistent data-collection environments in which meaningful consent, bystander privacy and the practical possibility of refusal are difficult to sustain \cite{MartinezMartin2021AmbientIntelligence,Cohen2020RemoteMonitoringDevices,WHO2021AIEthics}. Governance should begin with purpose limitation, data minimization and accessible controls, while treating passive non-use and temporary disengagement as legitimate patient choices.

Physical agency also makes safety time-dependent. An error can propagate from sensor drift to state estimation, planning, actuation and delayed recognition of harm. Risk management must therefore cover the full closed loop, including foreseeable misuse, human factors, degraded operation, safe-state transitions and post-deployment change, rather than certify isolated components once \cite{ISO14971,Yang2017MedicalRoboticsAutonomy}. Human oversight is necessary but not sufficient: a nominal ``human in the loop'' may have too little time, information or authority to intervene. Meaningful oversight requires calibrated autonomy, explicit action boundaries, interpretable escalation signals, logging, and tested handover procedures for operators with defined responsibility \cite{vanDeSande2026MeaningfulOversight}. High-risk actions should be constrained by independent safety monitors and conservative fallback modes, with validation in the populations, environments and workflows in which the system will actually act.

Fairness must likewise be evaluated at the level of the embodied pathway. Measurement bias can arise before inference---from skin-contact sensors, device placement, mobility differences, home infrastructure or caregiver support---and unequal access can determine who is represented in the data at all. Well-known failures of pulse oximetry and population-health algorithms show that apparently technical design choices can distribute missed detection and resources unequally \cite{Sjoding2020PulseOximetry,Obermeyer2019RacialBias}. Audits should therefore stratify not only model performance but also sensing quality, abstention, intervention burden, failure recovery and clinical benefit; they should examine intersections of demographic, physiologic and socioeconomic factors and include communities facing health-data poverty \cite{Liu2022MedicalAlgorithmicAudit,Chen2023AlgorithmFairness,Ibrahim2021HealthDataPoverty}. Because digital interventions can amplify existing service inequalities, evaluation must also measure who can obtain, maintain and safely use the embodiment \cite{Veinot2018InformaticsInequality}.

Finally, cybersecurity is inseparable from patient safety once software can affect physical state. The WannaCry disruption of the UK National Health Service illustrated how cyber incidents can interrupt care at system scale \cite{Ghafur2019WannaCry}. Embodied platforms require secure update mechanisms, authenticated commands, least-privilege interfaces, resilient local operation, vulnerability disclosure and a documented software bill of materials. Current medical-device guidance emphasizes cybersecurity across the product life cycle rather than as a premarket documentation exercise \cite{FDA2026Cybersecurity}. Accountability should be equally continuous: developers, deployers and clinical organizations need traceable logs, change control, incident reporting and a clear allocation of responsibility for monitoring, maintenance and override. The central governance principle is proportionality: the more a system can observe, decide and act without contemporaneous human control, the stronger the evidence, containment, transparency and institutional oversight it should require.

